\chapter{The Versioned Source Basis}\label{ch03:source}

\section{Why a source chapter exists}\label{ch03:why}

A monograph whose central discipline is that every claim must be traceable
to an exact artifact owes its reader, early and in one place, a complete
account of what its artifacts \emph{are}: which documents govern, how each
is identified, what authority each does and does not carry, and where the
read-only boundaries lie. This chapter is that account. Every later chapter
that states a fact about the engineering baseline, the audit, the
implementation, or the downstream formal kernels resolves its citation
through the identities fixed here; and where the sources do not specify a
fact --- a page count, a hash, a status --- the monograph's rule is to say
``unspecified'' rather than to estimate. The chapter also develops, as
methodology rather than housekeeping, the identification scheme itself:
why the work identifies its principal source by a source ID plus a
content hash instead of a public semantic version, and what a content
hash does and does not attest.

The reader should hold one asymmetry in mind throughout. The source basis
described here is \emph{versioned} in the strong sense --- every
controlling artifact is bound by a SHA-256 digest, and the publication
delta is bound to an exact repository commit --- but it is not thereby
\emph{validated} in any sense. Identity discipline establishes that all
parties are talking about the same bytes. It establishes nothing about
whether those bytes are correct, safe, complete, or authoritative; indeed,
two of the monograph's central results are demonstrations that the
principal source is internally defective at precisely identified
locations. Fixing identity is what makes such demonstrations meaningful:
an adverse finding against an unidentified document is an anecdote, while
an adverse finding against
\textsc{src-pcba-engineer-20260726} at a cited page anchor is a
reproducible check.

\section{The engineering baseline}\label{ch03:baseline}

\subsection{The artifact}\label{ch03:baseline:artifact}

The immediate engineering baseline of CAS PCBA 8 is a single large
document: the owner-supplied CAS PCBA Engineer build,
\texttt{CAS\_PCBA\_Engineer.pdf}, 743 pages, identified as

\begin{center}
\begin{tabular}{@{}ll@{}}
\toprule
Source ID & \textsc{src-pcba-engineer-20260726} \\
SHA-256 & \texttt{294320B39FA521EAC086E99E2D9DB913E2B6CFF0754FF82DB1EF966424566274} \\
Extent & 743 pages (PDF) \\
Role & Engineering reference baseline \\
\bottomrule
\end{tabular}
\end{center}

The document is an engineering reference: it contains the CAS circuit
description language (CCDL) exposition and its consolidated grammar, the
five published reference boards, the component and validator prose, the
appendix material, and the design narratives that the candidate
implementation and the source audit take as their object. Two properties
of the baseline shape everything downstream. First, it is \emph{large}:
at 743 pages, no reader or reviewer holds it whole in mind, so
page-anchored citation and mechanical audit are not conveniences but the
only available rigor. Second, it is \emph{closed as supplied}: the
candidate treats the build as a fixed input identified by its hash, and
does not amend, patch, or re-render it. Where the audit finds the
baseline internally inconsistent --- and it does --- the inconsistency
is preserved and formalized against the fixed bytes, never repaired in
place. A repaired design would be a new artifact with a new identity, a
rule stated in the audit and honored throughout this monograph.

\subsection{The dual-key identification}\label{ch03:baseline:dualkey}

The baseline is identified by a \emph{pair}: a human-legible source ID,
\textsc{src-pcba-engineer-20260726}, and a machine-checkable content
hash. The two keys do different work. The source ID is a stable name for
discourse --- it appears in claim records, review matrices, and prose,
and it encodes the supply date of the build (2026-07-26) for human
orientation. The hash is the actual identity: any party holding a file
alleged to be the baseline can compute its SHA-256 and compare all
sixty-four hex digits against the value above, and agreement identifies
the bytes with overwhelming confidence while any disagreement ---
substitution, truncation, re-rendering, silent edit --- is certain to
change the digest. The source ID without the hash would be a label
anyone could reuse for different bytes; the hash without the source ID
would be unreadable in prose. The pair is the identification scheme, and
the monograph never cites the baseline by either key alone in a
load-bearing context.

\subsection{The dual page-anchor convention}\label{ch03:baseline:anchors}

Citations into the baseline use a dual anchor, and the convention is
fixed here because every audit finding depends on it. The baseline, like
most large typeset books, has two pagination systems: the \emph{PDF
page}, counted from the first sheet of the file, and the \emph{book
page}, the printed folio, which lags the PDF count by the extent of the
front matter. The audit records both wherever both are meaningful ---
for example, the reserved Part-IV statement leaders appear at PDF
p.~578, book p.~562; the erratum language for missing appendix forms at
PDF pp.~591--592, book pp.~575--576; the consolidated structural EBNF at
PDF pp.~598--599, book pp.~582--583; the Chapter~18 safety-board listing
at PDF p.~342, book p.~326; its Appendix-D repetition at PDF p.~608,
book p.~592; and the Appendix-D note on the infinite-gain idealization
at PDF p.~610, book p.~594. The monograph's own citations use the PDF
anchor as primary --- it is the anchor a party holding the hashed file
can check without interpreting folios --- and repeat the book anchor
where the audit supplies it. Both anchors, again, are defined only
against the hashed build: a revised edition would renumber both systems
independently, which is one more reason the pin of
Section~\ref{ch03:baseline:dualkey} precedes everything else.

\section{Why not a semantic version}\label{ch03:version}

\subsection{The internal-label problem}\label{ch03:version:label}

The natural bureaucratic move for a 743-page evolving engineering book
would be to assign it a public semantic version --- ``v8.1,'' ``2026.3''
--- and cite that. The candidate deliberately refuses this move, for a
reason recorded in the Gate~4 audit: \emph{the baseline's own latest
internal changelog label predates the supplied 2026-07-26 build}. The
document's self-description of its version history, in other words, does
not extend to the state actually supplied; content exists in the build
that postdates the last internally recorded version label. Under that
condition, any public semantic version the candidate might attach would
be an invention --- a version fact fabricated by the citing party and
attributed, implicitly, to the source. There would then be two
authorities for the baseline's version (its own changelog and the
candidate's label), disagreeing silently, with the disagreement invisible
to every downstream reader.

The refusal generalizes into a rule the monograph applies to all source
facts:

\begin{ruleenv}[No invented source facts]\label{ch03:rule:noinvent}
The candidate does not assert, about a source, version facts, page
facts, count facts, or status facts that the source does not itself
carry. Where a needed fact is absent or stale in the source --- as the
version label is here --- the candidate substitutes an identification it
can itself warrant (a supply-dated source ID plus a computed content
hash) and discloses the substitution. Unspecified means unspecified.
\end{ruleenv}

\subsection{What is lost and what is gained}\label{ch03:version:tradeoff}

The cost of refusing a semantic version is real and admitted: semantic
versions carry ordering and compatibility conventions --- readers infer
from ``8.1 versus 8.2'' a direction of change and a rough promise about
its scope --- and a hash carries neither. Two hashes are merely
different; nothing about
\texttt{294320B3\ldots} orders it against any other build. The
compensation is that everything a hash \emph{does} say is true and
checkable, whereas the ordering a fabricated version would say is
unwarranted. For a research discipline whose failure mode is precisely
the silent promotion of unwarranted statements (Chapter~\ref{ch01:introduction}),
the trade is not close. Ordering among builds, where the work needs it,
is recovered honestly: by supply dates in source IDs (which claim only
chronology of supply, a fact the candidate can warrant) and by the bound
repository history, whose commits are themselves content-addressed and
totally ordered by their ancestry.

\section{Source-identity methodology: content addressing for large design documents}\label{ch03:methodology}

\subsection{The general practice}\label{ch03:methodology:practice}

Identifying an artifact by a cryptographic digest of its content ---
\emph{content addressing} --- is established practice wherever
computation must be accountable to exact inputs. Version-control systems
address every stored object by content hash, so that a commit identifier
fixes an entire tree of sources; package ecosystems pin dependencies by
digest so that a build consumes exactly the audited bytes; and the
reproducible-builds movement in software distribution (an effort across
compiler, packaging, and distribution communities, named here inline as
context rather than cited) pushed the discipline end to end, defining a
build as reproducible precisely when independent parties, from identified
inputs, produce bit-identical outputs --- with the comparison performed,
of course, by hash. The common structure in all three practices is the
one this chapter has already used: a digest converts ``the same
artifact'' from a social claim into a computation, and thereby converts
``I checked artifact X'' into a statement a third party can re-ground
without trusting the checker.

Applying the practice to a large \emph{design document} --- rather than
to code --- is less common and worth a methodological remark, because
documents are usually cited loosely (by title, edition, or version
label) on the assumption that editions are stable and public. The
baseline here is neither: it is a privately supplied, actively evolving
build whose internal version labeling has already been shown to lag its
content. In that setting the loose-citation assumption fails, and the
document must be treated the way a build input is treated: pinned by
digest, with every page-anchored citation in the audit and in this
monograph implicitly indexed by that digest. A page anchor such as
``PDF p.~342'' is meaningful \emph{only} against
\texttt{294320B3\ldots}; against any other build of the book, the anchor
is undefined. The audit's page anchors and the monograph's repetitions
of them inherit their precision entirely from the pin.

\subsection{Hash conservatism}\label{ch03:methodology:conservatism}

Content addressing is deliberately \emph{over}-sensitive, and the
over-sensitivity should be owned rather than hidden. A SHA-256 digest
changes under any byte change whatsoever, including changes with no
semantic content: re-rendering the same manuscript through a different
PDF producer, updating an embedded timestamp, recompressing an image
stream. Two files could, in principle, present identical pages to a
reader and carry different digests, and the scheme would treat them as
different artifacts. This is the right failure direction for the present
purpose. The alternative --- any scheme that tries to hash ``the
content'' rather than the bytes --- requires a canonicalization step
whose own correctness then silently joins the trusted base, and for a
743-page rendered document no such canonicalization is trustworthy. The
scheme therefore accepts false distinctions (semantically identical
builds counted as distinct) in exchange for excluding false
identifications (distinct builds counted as identical), because for an
audit trail the second error is fatal and the first is merely
inconvenient: a party holding a semantically identical but
differently-rendered build simply does not hold the audited artifact,
and must obtain it before checking anchors. Conservatism here is a
choice with a direction, and the direction is fail-closed.

The same conservatism explains why the basis pins \emph{artifacts}
rather than \emph{claims}: a digest over the audit document fixes every
finding, anchor, and status literal in it at once, so the monograph
never needs to enumerate which sentences of a source it relies on ---
reliance on any of them is reliance on the digest, and a revised audit
would be a visibly new artifact rather than a silently shifted
authority.

\subsection{What a hash attests --- and does not}\label{ch03:methodology:limits}

\begin{boundary}[Limits of content addressing]\label{ch03:bnd:hash}
A content hash attests identity of bytes and nothing else. Agreement
with a published digest establishes that a party holds the same
artifact; it does not establish that the artifact is correct, complete,
internally consistent, safe, authorized, or authoritative, and it does
not establish who authored it or with what review. In particular, no
self-hash --- no digest computed by an artifact's own machinery over
its own content --- confers evidence status, license status, or
authority on that artifact; a hash-bearing receipt is exactly as
authoritative as the process that produced it, which for locally
generated receipts is not at all. This monograph's claim discipline
(Section~\ref{ch01:discipline}) already quarantines authority-from-receipt
inferences; this boundary restates the quarantine at the level of the
identification scheme itself.
\end{boundary}

The point is not hypothetical hygiene. The project's history includes
historical AI-authored, self-hashed ``licensing receipts,'' and the
governing documents are explicit that these are not an authenticated
reviewer or legal authority. The identification scheme of this chapter
is therefore deliberately \emph{austere}: hashes bind discourse to
bytes; all authority questions are answered elsewhere, by the authority
matrix (Section~\ref{ch03:authority}) and the claim records, and the
answer at the bound state is uniformly negative.

\section{The baseline's own disclaimer and the Phase-0 authority matrix}\label{ch03:authority}

Identification says what the baseline is; authority says what it may be
used to support. On authority, three independent instruments agree, and
their agreement is a fact of the source basis rather than an
interpretation of this monograph.

First, \emph{the book's own disclaimer}: the baseline presents itself
as an engineering reference and disclaims, on its copyright page (PDF
p.~4), safety certification, guaranteed model behavior, and legal or
medical advice --- stating that a Cognitive Circuit Board ``does what
its schematic, parameters, and test report say it does --- no more,
and no less'' --- precisely the statuses a careless reader might
project onto a confident 743-page technical book. Second, \emph{the
Phase-0 authority matrix}: the
project's authority-mapping instrument, which assigns to each artifact
class the authorities it carries, denies to the baseline both
\emph{safety certification} and \emph{physical telemetry} --- the book
certifies no safety property of anything, and it reports no measurement
of any physical system. (The matrix is a Phase-0 process instrument; it
is not among the hash-identified artifacts of
Table~\ref{ch03:tab:artifacts}, and binding it there by digest is an
open documentation item.) Third, \emph{the candidate itself}: the Gate~4
audit and specification, in constituting the baseline as their
engineering input, restate both denials as conditions of use. The
alignment is deliberate: the baseline is admitted \emph{as design
text} --- as the object the grammar audit parses, the board listings
quote, and the known-fault fixture is drawn from --- and under no
instrument is it admitted as evidence of safety, measurement,
certification, or empirical validity.

It is worth pausing on why the denials must be \emph{institutional}
rather than merely sincere. A 743-page book with circuit diagrams,
validator prose, and a chapter titled around safety generates, by its
sheer weight and genre, a presumption of authority in casual readers
that no front-matter sentence reliably dispels; and the documented
history of this project includes exactly the failure the matrix guards
against --- generated artifacts (self-hashed receipts) whose form
imitated authority their process never had. An authority matrix answers
the genre problem structurally: authority becomes a typed property
assigned per artifact class by an instrument outside the artifact, so
the question ``may this document support a safety claim?'' has a
lookup, not a vibe, for an answer. For the baseline the lookup answers
no, twice --- no safety certification, no physical telemetry --- and
the candidate's conditions of use repeat the answer so that no chain of
citation, however long, can launder the book's confidence into a
warrant it never carried.

This three-way denial does load-bearing work later in the monograph.
The Cognitive Safety Board finding (previewed in
Section~\ref{ch01:sep:analogy} and developed in the audit chapters)
would be impossible to state cleanly if the baseline carried safety
authority: a certified design with an unreachable protective threshold
would be a certification failure, a different and graver genre of
result. Because the authority matrix already places the baseline
outside safety authority, the finding is what it should be: an adverse
\emph{source-quality} result about design text, preserved as a
known-fault fixture, with no certified system anywhere in the blast
radius --- and, equally, with no safety assurance anywhere in the
candidate, since the same matrix denies that authority to the
candidate's own artifacts.

\section{The controlling repository artifacts}\label{ch03:artifacts}

\subsection{The publication delta and its base commit}\label{ch03:artifacts:commit}

The monograph describes a candidate whose executable and evidentiary
content lives in a repository, and the publication lane is bound to an
exact state of it: the Project Coherence Paper~6 delta at base commit
\texttt{a1cd990109e7561dfd03949bf0e1653461c7ae44} (cited in short form
as \texttt{a1cd9901}). ``Delta'' has a precise meaning here: the
publication artifacts organize and cite the existing Gate~4 candidate;
they do not modify the Gate~4 specifications, registries, or manifests.
The implementation surface at the bound state is the
\texttt{src/cas\_pcba\_v8/} package, comprising 14 modules; its tested
behavior is summarized by the suite identities in the Gate~4 artifact
manifest and reported later in the monograph.

\subsection{The artifact set}\label{ch03:artifacts:set}

Seven artifacts control the source basis beyond the baseline itself.
Their identities and roles, as recorded in the governing seed and its
underlying publication delta, are:

\begin{longtable}{@{}p{5.2cm}p{7.6cm}@{}}
\caption{Controlling repository artifacts at base commit \texttt{a1cd9901}.}\label{ch03:tab:artifacts}\\
\toprule
Artifact and role & SHA-256 \\
\midrule
\endfirsthead
\toprule
Artifact and role & SHA-256 \\
\midrule
\endhead
\bottomrule
\endfoot
\texttt{CAS\_CORE\_V8\_CANDIDATE\_SPEC\_v0.1.md} --- semantic identity,
evidence, and research-boundary kernel &
\texttt{5FE6F91D005C7268D32C3943C2851138B\newline
30E2492C7D73877A04B30685CFCA1FC} \\[2pt]
\texttt{CAS\_PCBA\_V8\_SOURCE\_AUDIT\_AND\_SPEC\_v0.1.md} --- normative
Gate~4 engineering candidate: source audit, adverse findings, type
system, DRC contract, simulation contract &
\texttt{19707B9213040A3B55DB51132C30B90A6\newline
004D15A611EC759906335F2E41450F0} \\[2pt]
\texttt{PCBA\_CCDL\_CONSTRUCT\_COVERAGE\_v0.1.json} --- machine-readable
grammar/corpus adverse evidence: construct coverage, four direct
counterexamples, one bounded listing &
\texttt{5BB797EC37CF289B4CBE8DF1256F07425\newline
C3DA13AA8CF2DA03E47DFCE1CB46808} \\[2pt]
\texttt{GATE4\_REVIEW\_MATRIX\_v0.1.json} --- implemented/open decision
ledger: 53 review criteria with 20 open at the bound state &
\texttt{44DD99A847DA742C3EBEBC06D64452A68\newline
68421EE12129E789D1CC790DAE41B74} \\[2pt]
\texttt{GATE4\_ARTIFACT\_MANIFEST\_v0.1.json} --- exact artifact and
test-suite identities: 54 hash-bound files, three focused-suite test-ID
digests &
\texttt{49102F9BCBAA2345804F5EE3CA708FD6A\newline
84E83F8CF5ABB51E6612893A090D57B} \\[2pt]
TKS v8 core formal specification --- downstream finite-arrow boundary
(read-only) &
\texttt{F740BA05AB7ADB2DCDF55E271016D127E\newline
45073B015497DE8982D44A49F0EA2C4} \\[2pt]
Noetic Topos 2.0 specification --- downstream ambient/resolved evidence
boundary (read-only) &
\texttt{156A5D8C5320F36E81E3B46473BA8A038\newline
B35672DCD28839A404D96A3DB1A763A} \\
\end{longtable}

\subsection{The Gate-4 receipt as a source fact}\label{ch03:artifacts:receipt}

The artifact manifest's content is itself part of the source basis, and
its headline quantities are cited so often in later chapters that they
are fixed here once, with their status stated in the same breath. At the
bound commit, the integrated Gate~4 receipt (dated 2026-08-04 in the
underlying publication delta) records 54 hash-bound artifacts and three
green focused suites, whose suite-level test-ID digests are:

\begin{center}
\small
\begin{tabular}{@{}lrl@{}}
\toprule
Suite & Methods & Suite test-ID SHA-256 \\
\midrule
CAS Core v8 & 61 & \texttt{E2AF5014DE28DD0702211EF32A2757B99B869A13FAFBB0B8D1B18A50A9F8FEF5} \\
CAS PCBA v8 & 82 & \texttt{D5F776C5F15758CD390470B5A542CCA110784DC5EBB9807B6BE41B2C80453470} \\
CAS-Ohm v1 & 65 & \texttt{2E81A685961286D944894A22EB93345B97211B4433B6C71EAF628D84CD8EB1EF} \\
\bottomrule
\end{tabular}
\end{center}

The combined focused run comprises 208 test methods and 272
parameterized subtests; the receipt further records 11 claim records,
all Ring~4 and \textsc{not licensed}; 53 review criteria with 20 open;
four direct EBNF counterexamples; eight CAS-Ohm source conflicts; and an
integrated Gate~1 pass. One negative fact in the receipt is as
controlling as the positive ones and is repeated wherever the suites
are: the repository-wide collector was \emph{not} clean at the bound
receipt --- it reported four pre-existing package-path namespace
collisions after 1{,}200 collected items --- so the focused suites
cannot support a repository-wide verification claim, and no chapter of
this monograph makes one. All of these numbers are historical,
source-bound software receipts consumed by this monograph as source
facts; they are not experiments, and this publication lane adds no new
runs to them.

\subsection{A further read-only mathematical authority}\label{ch03:artifacts:paper3}

One additional identity belongs in the basis because a later chapter's
geometry discussion depends on knowing where its authority does
\emph{not} lie. The projector-geometry theorem statements referenced by
the candidate are governed solely by the Formal Mathematics Paper~3
record at its own commit,
\texttt{3c8491db6f5559cf090909d9de0efa03d924f53f}, with source-draft
SHA-256
\texttt{679E1119CB084D9A429F6FCC26374B18FF1FB6252D7374631975B4FF56BAEFC4}.
The engineering artifacts of Table~\ref{ch03:tab:artifacts} carry
read-only dependency summaries of that record, not the mathematics
itself; in particular, the registered normalized Gram-kernel result
requires common rank $d>0$, and the stronger CAS capacity transfer
remains, at that record's own status, a proof gap in review at Ring~4.
Nothing in this monograph re-derives, strengthens, or promotes any of
it; the identity is recorded here so that the later citation is, like
every other, a pointer to exact bytes.

Three remarks on the set as a whole. First, the artifacts form a
\emph{layered} basis: the Core~v8 specification supplies the semantic
and boundary kernel that everything else presupposes; the PCBA audit and
specification is the normative engineering document of the candidate
itself; the coverage record is the machine-readable half of the
grammar finding, deliberately separated from prose so that the adverse
evidence is directly consumable by tools; and the review matrix and
artifact manifest are the ledgers --- what was decided and what exists
--- against which ``at the bound state'' statements resolve. Second,
the set is \emph{closed for this monograph}: a claim that cannot be
grounded in the baseline, these seven artifacts, the governing research
paper, or the bound commit's receipts is not made. Third, the two
downstream specifications appear in the same table as the candidate's
own artifacts but with a different mode of use, developed next.

\section{The publication lineage of this monograph}\label{ch03:lineage}

The monograph itself has a documented lineage within the source basis,
and stating it closes the loop of self-description. The governing seed is
the research paper \emph{CAS PCBA 8: adaptive cognitive circuit
engineering with typed, fail-closed design rules and a preserved adverse
source audit}, version 0.2, dated 2026-08-07, which supersedes the white
paper v0.1; beneath it sits the Project Coherence charter Paper~6
engineering draft v0.1, source-bound to base commit
\texttt{a1cd990109e7561dfd03949bf0e1653461c7ae44}. The monograph expands
the seed to book length; it introduces no fact, hash, page anchor,
count, or status not present in the seed, the Paper~6 draft, or the
Gate~4 artifacts of Table~\ref{ch03:tab:artifacts}, and where the seed
is silent the monograph is silent.

The lineage's status fields are themselves source facts, recorded at
the Paper~6 draft and unchanged by anything in this text: authorship is
\textsc{open}, with no authorship or submission decision recorded;
internal cross-review is \textsc{open}; external specialist review is
\textsc{open}; and the submission state is \textsc{not ready for
submission}. The Paper~6 draft additionally binds its assertion
discipline to two machine-readable companions --- a claim-assertion
matrix and an evidence map, with citation identifiers resolving through
a citation ledger --- and declares that the publication lane modifies
no Gate~4 specification, registry, or manifest. The monograph inherits
all of it. In particular, the existence of this book-length treatment
changes no status: a longer exposition of an unlicensed, unreviewed,
Ring-4 candidate is a longer unlicensed, unreviewed, Ring-4 candidate.

\section{The audit's method, disclosed at the source level}\label{ch03:auditmethod}

The adverse findings that later chapters present rest on an audit whose
method is bounded, and the bounds are source-basis facts because they
determine what kind of evidence the findings are. As recorded in the
seed: text extraction over the hashed PDF located candidate pages;
every cited page was then rendered from the source PDF and visually
inspected, so that no anchor rests on extraction alone; and the
grammar comparison is between \emph{published surface syntax} and the
\emph{consolidated EBNF} --- it is not a parser implementation and not
a semantic-validation result. The machine-readable coverage record
(\texttt{PCBA\_CCDL\_CONSTRUCT\_COVERAGE\_v0.1.json},
Table~\ref{ch03:tab:artifacts}) fixes the outcome shape the method can
support: four boards are direct syntactic counterexamples to the
published grammar; one board is a bounded listing for which no direct
syntactic counterexample was established; and the corpus-level
conclusion is that the formal EBNF does not cover the published
five-board corpus.

Two consequences of the bounds are worth drawing here rather than in
the audit chapters, because they are consequences about \emph{evidence
class}. First, the method's verdicts are asymmetric in strength: a
counterexample is strong (an observed surface form with no producing
grammar alternative is a definite non-coverage witness at the cited
anchor), while the single non-counterexample verdict is weak by
construction (``no direct counterexample was established under a
bounded audit'' is not ``this board conforms''), and the source
documents are careful to claim only the weaker form --- the
corpus-level conformance claim fails \emph{without} overstating the
result for the one bounded listing. Second, the method's ceiling is
syntactic: even a hypothetical five-for-five conformance result would
have said nothing about semantic validity, and the audit accordingly
never uses the word. The candidate's own obligations mirror the
ceiling: what it demands of any future grammar --- an exact
construct-coverage matrix, rejection of unsupported forms, canonical
round-trips --- is precisely the machinery that would raise audits of
this kind from bounded surface comparison to checked conformance.

\section{Downstream read-only boundaries}\label{ch03:boundaries}

Two formal frameworks sit downstream of the candidate, and the source
basis fixes the candidate's relationship to each as strictly
\emph{read-only}: the candidate consumes their published boundaries and
conforms to them; it does not modify them, extend them, or acquire any
theorem from them by adjacency.

\subsection{TKS v8: the finite-arrow boundary}\label{ch03:boundaries:tks}

The TKS v8 core formal specification (hash in
Table~\ref{ch03:tab:artifacts}) defines the finite reference kernel
into which, under narrow conditions, a PCBA component's behavior may be
adapted. The boundary the candidate reads off that specification is the
\emph{finite-arrow} discipline: only a certified total, finite,
deterministic, stateless, non-stochastic component transition may be
represented as a finite arrow of the kernel; typed sequences become
semigroupoid paths; only an endomap has powers, so iterated feedback
requires a certified endomap. Everything richer --- continuous simplex
profiles, response curves, instrumented observations, stateful devices,
stochastic kernels, memristive history variables --- remains outside
the finite arrow and must be handled by distinct types and comparison
maps rather than coercion. The adapter that enforces this boundary, and
the immutable certificate it issues (binding the exact TKS kernel bytes
and profile, the board, component, and transition identities, the
tagged carrier, the semantic profile, and the fixed research boundary),
are described in their own chapter; what belongs to the source basis is
the direction of authority. The kernel's specification is an input the
candidate must satisfy, and satisfaction buys exactly one thing:
type-correct representation. The adapter carries no automatic CAS
semantic theorem, and no property of the kernel --- its finiteness,
its determinism, its formal pedigree --- transfers to the cognitive
content of any adapted component.

\subsection{Noetic Topos 2.0: the ambient/resolved evidence boundary}\label{ch03:boundaries:topos}

The Noetic Topos 2.0 specification (hash in
Table~\ref{ch03:tab:artifacts}) supplies the second read-only boundary,
this one about evidence. The distinction the candidate inherits is
between \emph{ambient} and \emph{resolved} standing for records: raw or
partially resolved records may be treated as external evidence
associated with the ambient construction, but a marked cover does not
make a measurement true, and records do not enter resolved semantics
until restriction and gluing under one calibration profile are defined
and tested. For the source basis, the inherited discipline lands in one
sentence: every observation-like record the candidate produces ---
test-point observations, measurement chains, trace events --- is
ambient-side material, and nothing in the candidate performs, or claims
to perform, the restriction-and-gluing obligations that would move any
record to resolved standing. The boundary is read off the downstream
specification and honored; it is not implemented, extended, or
discharged here.

\section{The Gate-5 preregistration kernel as read-only context}\label{ch03:gate5}

One further document sits in the source basis and requires the most
careful statement of all, because its entire function is to \emph{not}
be used yet. The Gate-5 preregistration specification is the project's
control kernel for future preregistered experimentation --- the
instrument through which, if the research program ever reaches
empirical work, calibration and held-out validation would be registered
in advance. Its status at the bound commit is:

\begin{center}
\begin{tabular}{@{}ll@{}}
\toprule
Instrument & Gate-5 preregistration specification \\
Status at bound commit & \textsc{open / not\_started / ring 4 / not licensed / held / not frozen} \\
Authority created & none (no live, scientific, canon, product, license, or disclosure authority) \\
Relation to this monograph & read-only context; supplies no experiment \\
\bottomrule
\end{tabular}
\end{center}

\begin{boundary}[Nothing here promotes Gate 5]\label{ch03:bnd:gate5}
The Gate-5 specification supplies no experiment to this monograph, and
this monograph performs no act with respect to it. Describing the
kernel, citing its status literals, or scheduling future preregistered
work through it in the research program does not start it, freeze it,
license it, or advance any of its status fields. Each literal in its
status line is load-bearing and is reproduced without softening:
\textsc{open} (the gate is not passed), \textsc{not\_started} (no
preregistered activity has begun), \textsc{ring 4} (outermost research
ring), \textsc{not licensed} (no claim licensing exists),
\textsc{held} (the artifact is held, not released), and
\textsc{not frozen} (the specification itself is not yet fixed ---
even the instrument's final form remains open). A reader who
encounters, anywhere downstream of this monograph, a suggestion that
Gate-5 work is underway or that its kernel validates anything in these
chapters is encountering a claim this source basis contradicts.
\end{boundary}

The kernel nonetheless belongs in the source basis rather than in a
footnote, for a structural reason. The research program with which the
monograph closes is dependency-ordered, and its empirical steps route
through Gate~5 by design: the candidate's position is that when
calibration and held-out validation happen, they must happen
preregistered, under a control instrument that existed --- with a known
identity and an honest \textsc{not\_started} status --- long before any
data did. Recording the kernel now, in its unstarted state, is what
makes that future discipline checkable: the preregistration instrument
is intended to be provenance-bound --- its digest is not yet recorded in
this volume's artifact table, and adding it is an open documentation
item, without which the guarantee that no later account can quietly
backdate an experimental plan into it rests on program practice rather
than on a binding checkable from this text. Read-only context today is the
commitment device for licensed evidence tomorrow; the asymmetry between
those two dates is the entire point.

\section{What the source basis excludes}\label{ch03:exclusions}

The discipline of Rule~\ref{ch03:rule:noinvent} has a complement: it is
equally a source-basis fact that certain things are \emph{not} in the
basis, and later chapters lean on these exclusions as often as on the
inclusions. Four exclusions bear listing.

First, \emph{no physical artifact}: the basis contains no fabricated
board, no instrument chain, no telemetry, and no physical measurement
of any kind; the word ``board'' throughout refers to a typed software
object or to a published design listing in the baseline. Second,
\emph{no live model or empirical dataset}: no calibrated live system,
no held-out evaluation data, and no experimental records exist in the
admitted basis; the only quantitative receipts are software-test
receipts bound in the Gate~4 manifest. Third, \emph{no external
review}: the basis records internal cross-review as open and external
specialist review (systems, safety, regulatory, clinical, formal) as
open; no artifact in Table~\ref{ch03:tab:artifacts} has passed
independent review, and the review matrix says so explicitly, with 20
of its 53 criteria open at the bound state. Fourth, \emph{no
authority}: as Sections~\ref{ch03:methodology:limits}
and~\ref{ch03:authority} establish from two directions, nothing in the
basis carries safety certification, measurement authority, claim
license, canon status, or disclosure authorization, and every claim
record at the bound state is Ring~4 and \textsc{not licensed}.

\section{Summary}\label{ch03:summary}

The source basis of this monograph is: one 743-page engineering
baseline, pinned as \textsc{src-pcba-engineer-20260726} with SHA-256
\texttt{294320B39FA521EAC086E99E2D9DB913E2B6CFF0754FF82DB1EF966424566274},
admitted as design text under a three-way denial of safety and
telemetry authority; seven controlling repository artifacts, each
hash-bound, layering a semantic kernel, a normative engineering
candidate, machine-readable adverse evidence, and two ledgers beneath
two read-only downstream boundaries; a publication delta at base commit
\texttt{a1cd9901} that organizes but does not modify the Gate~4 state;
and one deliberately unstarted preregistration kernel recorded in its
entirety of status literals precisely so that nothing can later be
smuggled into it. The identification scheme is content addressing with
disclosed limits: hashes make every citation in the chapters that
follow a checkable act, and confer authority on none of them. The
baseline is identified, not endorsed; the audit chapters will show, at
page-anchored resolution, that the distinction was necessary.
